\section{\textsl{base}}\label{sec:base}

\pbase{} is the entry-point of \prettytex{}. It is required by all other packages. It is used to set global options like language, page layout and such and loads elementary packages. 


\subsection{Package options}

\begin{table}[H]
    \centering
    \begin{tabular}{r|p{12cm}|l}
        \hline\textbf{option} & \textbf{description} & \textbf{type}\\\hline
        lang & the language to use throughout the document (mainly used for babel), expect valid babel language & \verb|string|\\
        singleLineCheck & the package \textsl{caption} provides the option singlelinecheck which allows to toggle whether or not captions of floats are centered (indented) or not via a boolean flag. \verb|false| disables centered captions, default is \verb|true| & \verb|boolean|\\
        colouredLinks & whether or not to use colouredLinks for reference, hyperrefs etc. default is \verb|true|. If \verb|false|, links will instead be highlighet by a coloured underline & \verb|boolean|\\
        parSkip & the \textsl{parskip} package introduces some small amount of vertical space between paragraphs instead of indenting the new paragraph. If \verb|true|, the pacakge is loaded. Default is \verb|false| & \verb|boolean|\\
        beamer & the \textsl{beamer} package clashes heavily with external page layouts, this if this option is passed, our page layout is disabled. Default is \verb|false| & \verb|boolean|
    \end{tabular}
    \caption{Package options for \pbase{}}
    \label{tab:base_options}
\end{table}


\subsection{Loaded packages}

\begin{table}[H]
    \centering
    \begin{tabular}{r|p{8cm}|p{4cm}}
        \hline\textbf{package} & \textbf{description} & \textbf{options}\\\hline
        \package{fontenc} & T1 font encoding ensures that non-ascii characters like the german umlaute (ö,ä etc) are actual glyphs instead of accented letters. This also improves copy-paste behavior & \option{T1}\\
        \package{inputenc} & allows for utf8 characters to bes used in source files (e.g. ä.ö. etc) & \option{utf8}\\
        \package{lmodern} & use latin modern font instead of computer modern\\
        \package{float} & allows for better control of float placement\\
        \package{geometry} & used for page layout (mostly for margins)\\
        \package{xcolor} & & \option{table}\\
        \package{cellspace} & & \option{column=C}\\
        \package{hyperref} & place clickable links in the text (if enabled via options, links are coloured instead of underlined) & \option{hypertextnames=false}\\
        \package{translations} & allows for automatic translations of text (e.g. theorem names)\\
        \package{contour} & give text an outline of any colour and width & \option{outline}\\
        \package{xspace}&\\
        \package{subcaption}&\\
        \package{graphicx}&\\
        \package{parksip} & only loaded if enabled via option, see \ref{tab:base_options} for explanation\\
        \package{enumitem} & more control over enums, specifically their labelling, only loaded if beamer is \verb|false| & \option{inline}
    \end{tabular}
    \caption{Packages loaded by \pbase{}}
\end{table}

\subsection{Document geometry}

The options passed to geometry are
\begin{multicols}{3}
    \begin{itemize}
        \item \option{a4paper},
        \item \option{left=2cm},
        \item \option{right=2cm},
        \item \option{top=1.25cm},
        \item \option{bottom=1.25cm},
        \item \option{includeheadfoot}
    \end{itemize}
\end{multicols}


\subsection{Colours}

\prettytex{} provides a few colours:

\begin{table}[H]
    \centering
    \begin{tabular}{r|l|lll |l}
        \hline\textbf{color-name}          & \textbf{defined as} & \textbf{R} & \textbf{G} & \textbf{B} & \textbf{hex-code}\\\hline
        \col{purple}{purple}               & RGB  & 180  & 2    & 212  & \#b402d4\\
        \col{green}{green}                 & RGB  & 0    & 153  & 51   & \#009933\\
        \col{red}{red}                     & RGB  & 204  & 0    & 0    & \#cc0000\\
        \col{blue}{blue}                   & RBG  & 0    & 102  & 255  & \#0066ff\\
        \col{orange}{orange}               & RGB  & 227  & 102  & 0    & \#e36600\\
        \col{darkred}{darkred}             & HTML &      &      &      & \#e05263\\
        \col{darkpurple}{darkpurple}       & HTML &      &      &      & \#5a016a\\
        \col{pastelpurple}{pastelpurple}   & HTML &      &      &      & \#a580d1\\
        \col{pastelpink}{pastelpink}  (pastelpink)     & HTML &      &      &      & \#f8c8dc\\
        \col{lightgrey}{lightgrey} (lightgrey)        & HTML &      &      &      & \#dddddd\\
        \col{darkpink}{darkpink}           & HTML &      &      &      & \#e6acac\\\hline
        \col{themecolor}{themecolor}       & HTML &      &      &      & \#00a8ff\\
        \col{themecolor2}{themecolor2}     & HTML &      &      &      & \#659157\\
        \col{themecolor3}{themecolor3}     & HTML &      &      &      & \#3d518c\\
        \col{themecolor4}{themecolor4}     & HTML &      &      &      & \#f25f5c\\\hline
        \col{lightgreybg}{lightgreybg}  (lightgreybg)   & rgb  & 0.97 & 0.97 & 0.97 & \\
        \col{headrulecolor}{headrulecolor} (headrulecolor) & rgb  & 0.8  & 0.8  & 0.8  & \\\hline
        \col{sigmared}{sigmared}           & HTML &      &      &      & \#af0f1e\\
        \col{sigmayellow}{sigmayellow}     & HTML &      &      &      & \#ffd500\\
        \col{discord}{discord}             & HTML &      &      &      & \#5865f2
    \end{tabular}
    \caption{Colours defined by \prettytex{}}
\end{table}


\subsection{Formatting}

We also provide some formatting macros:

\begin{table}[H]
    \centering
    \begin{tabular}{r|l}
        \hline\textbf{macro} & \textbf{expands to}\\\hline
        \prettytexMacro{HighlightColour} & pastelpurple\\
        \prettytexMacro{InfoColour}      & themecolor\\
        \prettytexMacro{WarnColour}      & orange\\
        \prettytexMacro{ErrorColour}     & red
    \end{tabular}
    \caption{\prettytex-macros}
\end{table}

\begin{table}[H]
    \centering
    \begin{tabular}{r|l|p{8cm}|l}
        \hline\textbf{macro} & \textbf{arguments} & \textbf{description} & \textbf{example}\\\hline
        \vma{ebf}       & \argu{text} & emphasised and bold-faced text             & \ebf{example}\\
        \vma{ibf}       & \argu{text} & italicised and bold-faced text             & \ibf{example}\\
        \vma{ibf}       & \argu{text} & underlined and bold-faced text             & \ubf{example}\\
        \vma{eubf}      & \argu{text} & emphasised, underlined and bold-faced text & \eubf{example}\\
        \vma{iubf}      & \argu{text} & italicised, underlined and bold-faced text & \iubf{example}\\\hline
        \vma{col}       & \argu{colour}, \argu{text} & write text in colour         & \col{green}{example}\\
        \vma{REFL}      & \argu{text} & highlight reference that does not yet exist                & \REFL{example}\\
        \vma{highlight} & \argu{text} & highlight text      & \highlight{example}\\
        \vma{info}      & \argu{text} & info text      & \info{example}\\
        \vma{warn}      & \argu{text} & warn text      & \warn{example}\\
        \vma{error}     & \argu{text} & error text      & \error{example}\\
    \end{tabular}
    \caption{Formatting macros}
\end{table}

\begin{table}[H]
    \centering
    \begin{tabular}{r|l|p{6cm}|l}
        \hline\textbf{macro} & \textbf{arguments}               & \textbf{description} & \textbf{example/output}\\\hline
        \vma{cw}                   & \argu{text} or \argu{expression} & give the passed argument a white outline with width .06em & \cw{example}\\
        \vma{th}                   &                                  & write th as superscript (for dates) & 4\th\\
        \vma{prettytex}            &                                  & write prettytex in pacakge format & \prettytex
    \end{tabular}
    \caption{Miscellaneous macros}
    \label{tab:base_misc}
\end{table}

\subsection{Miscellaneous}

\begin{table}[H]
    \centering
    \begin{tabular}{r|l|p{8cm}|l}
        \hline\textbf{macro} & \textbf{arguments} & \textbf{description} & \textbf{example/output}\\
        \vma{tightlist} & & remove vertical spacing from all lists\\
        \vma{rom} & \argu{positive integer} & convert input to lowercase roman numeral & \rom{42}\\
        \vma{Rom} & \argu{positive integer} & convert input to uppercase roman numeral & \Rom{69}\\
        \vma{alp} & \argu{integer in [1,26]} & convert input to lowercase alph-value    & \alp{24}\\
        \vma{Alp} & \argu{integer in [1,26]} & convert input to uppercase alph-value    & \Alp{16}\\
        \vma{listfigures}  & & this is a verbatim copy of the built-in macro \vma{listfigures} but allows the list to be included as a subsection &\\
        \vma{listtables}   & & this is a verbatim copy of the built-in macro \vma{listtables} but allows the list to be included as a subsection &\\
        \vma{autoheadings} & & set header and footer to sensible default based on document metadata &\\
    \end{tabular}
    \caption{Miscellaneous macros from \pbase}
\end{table}